\subsection{\bt{Optimization revisions}{优化修订}}
\label{sec:opt-eval}
\bi{\emph{Setup (Q7).} We add one optimization sweep to the end-to-end protocol of Section~\ref{sec:scenarios}, after learning and before the held-out questions, at one million sales rows, and run it once. \emph{Result.} \OptPublished\ of the \OptDefs\ published definitions of the five metrics received a faster revision, each from the key-range rule; since rule candidates were published first, no LLM proposal was needed. Paired execution time fell by \OptSaveMin--\OptSaveMax\% (mean \OptSaveMean\%), every 95\% interval below zero, with equal results on \OptPeriods\ periods per definition and as of learning time. Testing a published candidate took \AmOptTestMinS--\AmOptTestMaxS\,s (\AmOptTestTotalS\,s for all of them) and saves \AmOptSaveMinMs--\AmOptSaveMaxMs\,ms per execution, so each revision repays its test after \AmOptPaybackMin--\AmOptPaybackMax\ executions, counted over all consumers. Consumers adopted the new form: \OptHoldAdopted\ of the \OptHoldTasks\ held-out queries used the key-range predicate of the revision they retrieved. Over the held-out questions and the eleven changes, agents answer \OptAcc\ of \OptTasks\ questions, against \OptBaseAcc\ of \OptBaseTasks\ in a run of the same configuration without the sweep; the per-change differences are within the variation between runs. Under late-arriving date keys, the date-key completeness condition invalidates the key-range revisions exactly as it invalidates the join form, so an optimized revision keeps the protection of its predecessor.}
{\emph{设置（Q7）。}我们在第~\ref{sec:scenarios} 节的端到端流程中加入一轮优化，位于学习之后、留出题之前，数据为 100 万行门店销售，运行一次。\emph{结果。}五个指标的 \OptDefs\ 条已发布定义中有 \OptPublished\ 条得到更快的修订，均来自日期键范围规则；由于规则候选先发布，没有用到大模型的提议。配对执行时间下降 \OptSaveMin\%--\OptSaveMax\%（平均 \OptSaveMean\%），95\% 区间全部低于 0，每条定义在 \OptPeriods\ 个期间上以及按学习时刻结果相同。测试一个发布的候选用时 \AmOptTestMinS--\AmOptTestMaxS\ 秒（全部合计 \AmOptTestTotalS\ 秒），每次执行节省 \AmOptSaveMinMs--\AmOptSaveMaxMs\ 毫秒，因此每个修订在被所有使用者累计执行 \AmOptPaybackMin--\AmOptPaybackMax\ 次后收回其测试代价。使用者采用了新写法：\OptHoldTasks\ 条留出题查询中有 \OptHoldAdopted\ 条使用了所检索修订中的日期键范围谓词。在留出题与 11 种变化上，智能体答对 \OptTasks\ 道中的 \OptAcc\ 道，同一配置不做优化的一次运行为 \OptBaseTasks\ 道中的 \OptBaseAcc\ 道，各变化上的差异在运行间的波动范围内。在迟到日期键下，日期键完整性条件使日期键范围修订失效，与它使连接写法失效完全一样，因此优化后的修订保留了前一修订所受的保护。}
